% Revised online appendices for dSABRE.
% This file is intended to replace appendices.tex in the main manuscript.
% It assumes the main manuscript defines \dSABRE, \TeleSABRE, \dintra,
% \dcore, \dphys, and uses algorithm/algorithmicx and booktabs.

\appendices

\noindent\textbf{Scope.}
These appendices provide material needed to verify the algorithmic claims and
interpret the comparisons in the main paper.  They retain: (i) executable
routing semantics, (ii) the completion argument, (iii) the complexity
argument and its implementation validation, and (iv) the controlled
cross-model and architecture audits that support claims made in
Sections~III--V.  Exploratory experiments, superseded configurations, and
implementation-debugging history are intentionally omitted.  Complete
per-seed JSON records, commands, and patched comparison-tool environments are
released with the artefact repository.

% -----------------------------------------------------------------------------
\section{Completion Under Finite Capacity}
\label{app:safemode}

Theorem~\ref{thm:completion} is proved in full in
Section~\ref{sec:recovery}.  This appendix records the implementation
properties that proof takes for granted, the recovery cost it leaves out, and
how often the mechanism actually fires.

Let $f_c$ be the number of free physical slots in core $c$ and
$F_{\rm free}=\sum_c f_c=P-n$.  Both local SWAPs and teleports conserve
$F_{\rm free}$.  The router checks before routing that
\begin{equation}
 F_{\rm free}\ge 2K+1,
 \label{eq:safe_feasible}
\end{equation}
and that every core and the core graph are connected, each core has at least
five sites, and checkpoint--rollback is enabled.  An ordinary teleport is
legal only if its landing core has $f_{c_{\rm next}}\ge\phi_1=2$ before the
move.  It therefore leaves that core non-full; local SWAPs do not alter any
core occupancy.  Recovery additionally requires the per-core reserve
\begin{equation}
 f_c \ge 2 \quad\text{for every core } c,
 \label{eq:safe_invariant}
\end{equation}
which ordinary routing does not maintain: measured over every teleport of the
64-qubit suite, some core sits below it at $24.6\%$ of action boundaries
($11\,177$ of $45\,527$), reaching $46.4\%$ on AE---and at none of them does
any core reach zero, which is the floor actually enforced.

The properties the proof assumes are: local SWAPs and teleport
macro-actions are atomic, a failure restoring layout, inverse layout, remaining
DAG, counters and trace to the entry state; an ordinary teleport is generated
only when its landing core has $\phi_1=2$ free sites, and its implementation
uses no auxiliary teleport that changes a third core's occupancy; and the
initial layout obeys Eq.~\ref{eq:safe_invariant}, which
Section~\ref{sec:initial_mapping} enforces by reserving sites per core.  That
last one is not cosmetic: the global condition Eq.~\ref{eq:safe_feasible}
alone does not make an arbitrary initial placement recoverable, since every
free site could lie outside one completely full core, whose occupied source
port would then not be evacuable.

\paragraph*{Recovery cost}
Restoring the reserve moves at most $O(K)$ vacancies across paths of length at
most $D_K$, for $O(KD_K)$ teleports.  Co-locating the operands and restoring
the occupancy vector adds $O(D_K)$ teleports.  A recovery transaction therefore
uses $O(KD_K)$ EPR pairs in the worst case.  This is a feasibility bound, not
an EPR-optimality claim.

Three further facts follow from the same construction.  When the meeting core
already has headroom the transaction costs exactly $d$ teleports for a gate at
core distance $d$---the information-theoretic lower bound, since each teleport
advances one qubit by one core---and when it does not, the relay adds at most
$d$ more.  The $O(KD_K)$ bound is $\Theta(K^2)$ in the worst case, on a core
path whose slack all sits at one end.  And charging each gate its stall window,
its own transaction and one progressing iteration bounds the whole loop at
$|C|(L_\mathrm{deadlock}{+}2)$ iterations.



\begin{table}[t]
\centering
\caption{The shipped default---\dSABRE{} with the capacity
conditions, whose column reproduces Table~\ref{tab:main}'s 64-qubit EPR
figures---against the score-only design it replaced, over three initial
layouts $\times$ forward--backward--forward.  EPR pairs are the reported
route's; lower is better.  The mechanism counters below are totals over all
nine routes per circuit, on the \emph{search} convention of
Table~\ref{tab:tier2}, so they exceed what any single reported route shows.
A dash marks a mechanism the design does not have; a zero marks one it has and
never uses, and the distinction matters for rollback, which both designs
carry and neither ever takes.}
\label{tab:safemode}
\small
\setlength{\tabcolsep}{5pt}
\begin{tabular}{@{}lrrr@{}}
\toprule
Circuit & Score-only & Default & $\Delta$ \\
\midrule
ae          & $200$   & $165$   & $-17.5\%$ \\
ghz         & $6$     & $6$     & $0.0\%$   \\
graphstate  & $15$    & $15$    & $0.0\%$   \\
qft         & $176$   & $180$   & $+2.3\%$  \\
qnn         & $495$   & $539$   & $+8.9\%$  \\
random      & $718$   & $732$   & $+1.9\%$  \\
qpeexact    & $184$   & $187$   & $+1.6\%$  \\
qaoa        & $538$   & $555$   & $+3.2\%$  \\
multiplier  & $1285$  & $1285$  & $0.0\%$   \\
\midrule
Total       & $3617$  & $3664$  & $+1.3\%$  \\
Geo.\ mean  & ---     & ---     & $+0.2\%$  \\
\midrule
\multicolumn{4}{@{}l}{\emph{Mechanism counters, whole suite, all
  $81$ routes}} \\
Recovery invocations         & $111$   & $42$    & \\
\quad guaranteed transactions & ---    & $42$    & \\
\quad\quad of which failed   & ---     & $0$     & \\
Transaction rollbacks taken  & $0$     & $0$     & \\
Reactive room-making calls   & $104$   & $0$     & \\
Routing iterations           & $207\,650$ & $201\,776$ & $-2.8\%$ \\
\bottomrule
\end{tabular}
\end{table}

% -----------------------------------------------------------------------------
\section{Complexity Validation}
\label{app:complexity}

Section~III-H derives Eq.~\ref{eq:complexity} and is not repeated here.  This
appendix records the four structural properties the derivation relies on but
does not argue, and the measurements that check the derivation against the
implementation that produced every reported result.

The properties are: (1) the circuit DAG has bounded degree; (2) the front
layer is a matching, so a candidate SWAP changes at most two front-layer
distance terms; (3) BFS layers of the extended set are matchings, which is
what bounds the per-update cost of building $E$ at $O(n+L)$ after a gate
retirement and $O(L)$ for its per-core restrictions; and (4) gates retire only
from the front layer, so remaining in-degrees and rollback logs can be
maintained incrementally rather than rebuilt.  Property~(4) is the one that
would fail silently if it were violated, which is why it is measured below
alongside~(2).

\begin{table}[t]
\centering
\caption{Implementation checks supporting the complexity argument.}
\label{tab:appendix-complexity-checks}
\small
\resizebox{\ifdim\width>\columnwidth \columnwidth\else\width\fi}{!}{%
\begin{tabular}{@{}llr@{}}
\toprule
Claim & Measurement & Result \\
\midrule
Scaling at fixed architecture & fit versus CX count, 64q suite (9) & $N^{0.97}$ \\
Scaling with the machine & fit versus CX count, scalability (6) & $N^{2.2}$ \\
Front-layer matching & violations across four suites & 0 \\
Parallel local routing is rare & active cores per local iteration & 1.00--1.01 \\
Shared $E$ versus submitted sweep & total compile time, 39 circuits & $1745$ s $\to239$ s \\
Shared $E$ versus submitted sweep & geometric-mean EPR & $-2.4\%$ \\
64q shared $E$ versus sweep & compile-time improvement & $9.7\times$ \\
Hierarchical $\dphys$ & disagreeing pairs, seven architectures & 0 of 385,336 \\
\bottomrule
\end{tabular}%
}
\end{table}

Both scaling rows are ordinary least squares on $\log(\text{time})$
against $\log(\text{CX count})$, over \dSABRE{}'s per-seed compile time with
its default parameters; the fitted exponents are $0.97$ ($R^2=0.95$) on the
nine 64-qubit circuits and $2.17$ ($R^2=0.99$) on the six QFT/QPEexact
scalability instances, and the fixed-architecture exponent is $0.92$
($R^2=1.00$) and $1.21$ ($R^2=0.98$) on the 36- and 25-qubit suites.  The two
must be reported separately: within one suite only the circuit grows, whereas
the scalability series grows the device with it (six, twelve and twenty
$5{\times}5$ cores), which Eq.~\ref{eq:complexity} charges through $K^2$.  An
earlier draft quoted a single $N^{1.05}$ fit pooled over all 27 circuits; that
figure mixed the two regimes and, because the suite times were protocol totals
over three seeds while the scalability times were single routes, two timing
conventions as well.  The shared-$E$ comparison in
Table~\ref{tab:appendix-complexity-checks} is the relevant validation of the
implementation actually used for every reported result.  Earlier per-core
``taint'' sweeps, profiling detail, and alternative lookahead constructions
are omitted because they are not part of the algorithm evaluated in the main
paper.

\paragraph*{Timing protocol against the submitted version}
The submitted manuscript reported $801$\,s for a 360-qubit QFT and this one
reports $48$\,s, and the two are not measurements of the same thing.  The
architecture differs: the submitted figure was taken on an H-grid of six
$9{\times}9$ cores (486 physical qubits), while the scalability series
specified in Section~\ref{sec:largecircuits} holds core size at $5{\times}5$
and grows the core count, so 360 qubits is twenty cores and 500 physical
sites.  Two implementation changes account for the rest: the incremental
bookkeeping of 2026-08-07, which is output-identical to what it replaced and
was verified as such over 207 routing passes, and the replacement of the
per-core extended-set sweep by one shared construction, which is $43\times$
faster at 360 qubits.  The timing boundary is unchanged in both versions: the
reported number covers layout generation and every route of the
\texttt{SabreLayout}$\times$forward/backward/forward protocol, on the same
machine class and the same three seeds.

% -----------------------------------------------------------------------------
\section{Supplementary Empirical Audits}
\label{app:audits}

\subsection{Parameter robustness}
\label{app:sensitivity}
On the 64-qubit suite, every default is the minimum in its one-at-a-time
sweep: $w_e=0.25$, $|E|=20$, $\gamma=0.9$, and $c_{\rm pen}=15$.  Removing
lookahead ($w_e=0$) costs 41.4\% geometric-mean EPR; setting $|E|=5$ costs
14.9\%, $\gamma=0.5$ costs 20.7\%, and $c_{\rm pen}=0$ costs 9.4\%.
At 25 qubits the curves are largely flat except for removing lookahead, which
costs 11.4\%.  These results support the fixed defaults but do not establish
a universal parameter setting.  Table~\ref{tab:sensitivity} sweeps the cost
model instead of the score: \dSABRE{}'s margin widens monotonically as
$c_{\rm tele}$ rises from 10 to 100 at fixed $c_{\rm swap}=3$.

\begin{table}[t]
\centering
\caption{Cost-model sensitivity: geometric-mean cost reduction of
  \dSABRE{} over \TeleSABRE{} as $c_{\mathrm{tele}}$ varies
  ($c_{\mathrm{swap}}=3$ fixed), combined cost being
  $c_{\mathrm{tele}}\cdot\text{EPR}+c_{\mathrm{swap}}\cdot\text{SWAP}$.
  Negative is better.  Each suite is taken over its full circuit set:
  six at 25q and 36q, and at 64q the eight of nine on which \TeleSABRE{}
  converges---the same basis as Table~\ref{tab:main}'s $-49.1\%$.  The
  teleport cost cancels out of the routing decision itself
  (Section~\ref{sec:principle}), so this reweights the EPR/SWAP counts
  already in Table~\ref{tab:main} rather than re-routing; confirmed by a
  direct re-route sweep at $c_{\mathrm{tele}}\in\{10,100\}$, which
  reproduced identical EPR and SWAP counts.}
\label{tab:sensitivity}
\setlength{\tabcolsep}{6pt}
\small
\begin{tabular}{@{}r rrr@{}}
\toprule
$c_{\mathrm{tele}}$ & 25q & 36q & 64q \\
\midrule
10 & $\mathbf{-46.9}$ & $\mathbf{-55.8}$ & $\mathbf{-37.9}$ \\
20 & $\mathbf{-47.9}$ & $\mathbf{-57.6}$ & $\mathbf{-41.0}$ \\
50 & $\mathbf{-49.1}$ & $\mathbf{-59.6}$ & $\mathbf{-44.7}$ \\
100 & $\mathbf{-49.8}$ & $\mathbf{-60.7}$ & $\mathbf{-46.6}$ \\
\bottomrule
\end{tabular}
\end{table}

\subsection{Mechanism ablation at 25 qubits}
\label{app:mech25}
Table~\ref{tab:mech25} repeats the main paper's Table~IV at 25 qubits.  The
two capacity mechanisms remain substitutes at this scale, and the capacity
score alone is slightly negative ($-2.5\%$), which is why the main-paper row
is read together with the legality floor rather than on its own.

\begin{table}[t]
\centering
\caption{Ablation on the 25-qubit suite (B-grid $2{\times}2$ of
  $4{\times}4$, six circuits), the sparsely-loaded counterpart to
  Table~\ref{tab:mech}.  \textbf{Full} is Table~\ref{tab:main}'s own geometric
  mean.  Every row completes every circuit.}
\label{tab:mech25}
\setlength{\tabcolsep}{4pt}
\small
\begin{tabular}{@{}l rr c@{}}
\toprule
Configuration & gmEPR & $\Delta$ (\%) & aborts \\
\midrule
\textbf{Full} (baseline)                    & 15.3 & ---      & 0 \\
\midrule
\multicolumn{4}{@{}l}{\emph{Scoring terms}} \\
No capacity score term ($c_{\mathrm{pen}}{=}0$) & 14.9 & $-2.5$   & 0 \\
No extended-set lookahead ($w_e{=}0$)       & 17.1 & $+11.4$  & 0 \\
Topological extended set (not BFS)          & 15.3 & $+0.2$   & 0 \\
\midrule
\multicolumn{4}{@{}l}{\emph{Capacity as a mechanism}} \\
No legality floor (score term only)         & 15.1 & $-1.1$   & 0 \\
\textbf{Neither floor nor score term}       & 25.7 & $\mathbf{+67.7}$ & 0 \\
Strict floor ($\phi_1{=}r{+}1$)             & 15.4 & $+0.4$   & 0 \\
\bottomrule
\end{tabular}
\end{table}

\subsection{Topology sensitivity}
\label{app:corenet}
Table~\ref{tab:appendix-heavyhex} provides the per-suite results underlying
Section~IV-C.  The ring and star use four 27-qubit heavy-hex cores (108
physical qubits), rather than the six 16-site grid cores used for the 64q
H-grid.  The same circuit set, hyperparameters, and best-of-three-seed
protocol are used.

\begin{table}[t]
\centering
\caption{Geometric-mean EPR counts on the 64-qubit suite for non-grid core
networks.  A dash denotes a \TeleSABRE{} non-convergence.}
\label{tab:appendix-heavyhex}
\small
\resizebox{\ifdim\width>\columnwidth \columnwidth\else\width\fi}{!}{%
\begin{tabular}{@{}lrrr@{}}
\toprule
Architecture & \TeleSABRE{} & \dSABRE{} & $\Delta$EPR \\
\midrule
Four-core heavy-hex ring (9 circuits) & 177.4 & 85.2 & $-52.0\%$ \\
Four-core heavy-hex star (8 common circuits) & 142.4 & 83.3 & $-41.5\%$ \\
Four-core heavy-hex star (all 9, \dSABRE{} only) & --- & 101.1 & --- \\
\bottomrule
\end{tabular}%
}
\end{table}

For the 200-qubit QFT core-granularity control, six $7\times7$ cores require
353 EPR pairs versus 579 for \TeleSABRE{}; twelve $5\times5$ cores require
766 versus 1232.  The respective reductions are 39.0\% and 37.8\%.  This
single control indicates that the reported routing advantage is not confined
to one core granularity; it is not a general hardware co-design conclusion.

\subsection{TeleSABRE action audit}
\label{app:telesabre}
The direct comparison is with \TeleSABRE{} because both account for local
SWAPs and use one-EPR teledata moves.  The action spaces nevertheless differ.
\TeleSABRE{} generates a candidate only when the operand is already adjacent
to an empty source port and the selected destination port is empty.  \dSABRE{}
prices staging and port eviction in $d_{\rm prep}$ and enumerates each legal
neighbouring landing core.  Thus, availability is a scored macro-action cost
in \dSABRE{} rather than a condition that removes the action.

In a streamed decision trace, \TeleSABRE{} issued 348 of 414 communication
operations (84.1\%) while same-core front-layer work was pending on the
25-qubit suite; the corresponding count was 2368 of 3982 (59.5\%) on the
64-qubit suite.  This does not show that \TeleSABRE{} destroys local work:
none of 407 (25q) or 3725 (64q) teleports moved an operand of a same-core
front gate.  It establishes the narrower claim used in the paper: the routers
apply different ordering policies.

On the 64-qubit suite, telegate accounts for 289 of \TeleSABRE{}'s 4478 EPR
pairs (6.5\%) among converged circuits.  Its use is concentrated in AE
(26.6\%) and QAOA (13.8\%); it is at most 3.2\% for the other six converged
circuits.  The comparison should therefore be read as a controlled teledata
comparison with a small, reported telegate asymmetry, not as a claim that the
two action sets are identical.

\subsection{Telegate as a scored macro-action}
\label{app:telegate}
Section~VI states that scoring a single telegate is straightforward and
that grouping is the open problem.  We implemented the scored action to check
that claim rather than assert it.  For a front-layer gate whose operands sit in
\emph{adjacent} cores, one macro-action evicts both endpoints of a link, stages
each operand beside its own endpoint, consumes one EPR pair to cat-entangle one
operand onto the far endpoint, executes the gate there, and disentangles.
Neither operand changes core and the cat qubit is measured out inside the same
action, so both cores end at their entry occupancy: this is what licenses
$c_{\mathrm{cap}}=0$, and it also means the action cannot violate the capacity
invariant of Appendix~\ref{app:safemode}.  Restricting it to adjacent cores
keeps one pair live for one gate; a $k$-hop telegate would need $k$ pairs held
simultaneously, which is the lifetime assumption
Appendix~\ref{app:dmax} prices.  Candidates enter the same sorted list as
teleports and the best-scoring action of either kind is executed.  A telegate
receives no lookahead credit, because it moves no qubit between cores; that
asymmetry is the amortisation trade-off, not an omission.

Across seventeen variants of the score (flat bias, an explicit
amortisation penalty, a selective rule that offers a telegate only where
relocation would gain nothing in $E$, and weights on each term), none improves
on teledata alone.  The score as described above costs $+93.5\%$ EPR in
geometric mean over the nine 64-qubit circuits; the best variant reaches
$+1.0\%$, and only at settings where the telegate has almost stopped firing.
Table~\ref{tab:telegate_amort} gives the reason, measured from \dSABRE{}'s own
operation trace: a \emph{residence} is the interval a logical qubit spends in
one core between teleports, so a teleport-opened residence costs exactly one
EPR pair and serves however many two-qubit gates that qubit executes there.
One pair serves $15.75$ gates on this suite and $332$ in the worst case, while
an ungrouped telegate serves one by construction.  Telegate-whenever-legal
costs $6.1\times$ the baseline's EPR.  The action wins only where amortisation
is impossible: on capacity-tight variants of the same chip, GHZ ($0.83$
gates/pair) improves from $33$ to $16$ pairs and GraphState ($1.00$) from $69$
to $31$ --- both $64$-gate circuits in which each qubit interacts once, and
both cases a grouped telegate would not help either, since there are no
compatible groups to form.  Every routing reported here was checked by
replaying its emitted operation stream against the architecture and confirming
that the executed two-qubit gates are a valid topological order of the input
circuit.

\begin{table}[t]
\centering
\caption{What one teledata EPR pair buys on the 64-qubit suite (H-grid
  $2{\times}3$ of $4{\times}4$), from \dSABRE{}'s own trace.  \textbf{g/pair}
  is gates executed per teleport-opened residence; the last three columns
  partition those residences.  A single-gate telegate always serves exactly
  one gate, so it can only compete on the transit and one-gate columns.}
\label{tab:telegate_amort}
\setlength{\tabcolsep}{4pt}
\small
\begin{tabular}{@{}l rr rrr@{}}
\toprule
Circuit & EPR & g/pair & transit (\%) & one (\%) & $\geq 2$ (\%) \\
\midrule
AE          & 165  & 18.39 & 3.0  & 0.0   & 97.0 \\
GHZ         & 6    & 0.83  & 16.7 & 83.3  & 0.0  \\
GraphState  & 15   & 1.00  & 0.0  & 100.0 & 0.0  \\
QFT         & 180  & 17.49 & 3.3  & 0.0   & 96.7 \\
QNN         & 539  & 27.43 & 0.9  & 0.7   & 98.3 \\
Random      & 732  & 3.97  & 20.2 & 9.8   & 69.9 \\
QPE-exact   & 187  & 18.16 & 5.3  & 0.5   & 94.1 \\
QAOA        & 555  & 13.26 & 1.6  & 0.0   & 98.4 \\
Multiplier  & 1285 & 17.94 & 3.3  & 22.3  & 74.4 \\
\midrule
\textbf{Suite} & \textbf{3664} & \textbf{15.75} & 6.2 & 10.5 & \textbf{83.4} \\
\bottomrule
\end{tabular}
\end{table}

\subsection{Compile time}
\label{app:timing}
Table~\ref{tab:timing} gives the per-circuit compile times behind Section~IV-C.

\begin{table}[t]
\centering
\caption{Compilation time in seconds on the 64-qubit suite (H-grid
  $2{\times}3$ of $4{\times}4$).  Each column is the cost of the protocol that
  produced that tool's reported result, not of one attempt: \TeleSABRE{} needs
  a single seed, being deterministic here---its three seeds return identical
  EPR counts on all nine circuits---while \dSABRE{} needs three
  \texttt{SabreLayout} seeds, whose results genuinely differ (AE:
  $165/196/212$), and \texttt{pytket-dqc} five by its own convention.  The
  \TeleSABRE{} and \dSABRE{} columns were measured in isolation on one machine;
  the \texttt{pytket-dqc} column is the best-of-5 search behind
  Table~\ref{tab:fair}, taken from a separate run.  Timings of
  identical work vary by up to ${\sim}2.7\times$ with machine load, so the
  table should be read for orders of magnitude.  $^{\ast}$\TeleSABRE{} fails to
  converge on Random.  $^{\dagger}$\texttt{pytket-dqc} fell back to
  \textsc{PartitioningHeterogeneous} because the embedding distributor returned
  no distribution---it exhausts its budget on QNN and Multiplier, and
  raises on Random; the time spans the whole search, so it includes that
  abandoned attempt.  $^{\ddagger}$the embedding distributor did return on
  QAOA, but only two of its five seeds completed inside the per-circuit
  budget, so its time is set by that budget rather than by convergence.  The geometric mean is over the eight
  circuits all three compilers complete---the same basis in every column, so
  the columns are comparable---and excludes Random, which only \TeleSABRE{}
  fails.  The lower block is the QFT scalability series, whose
  architectures grow with the circuit ($6$, $12$ and $20$ cores of
  $5{\times}5$) and which is therefore not comparable across rows with the
  fixed-architecture suite above; at 360 qubits \TeleSABRE{} spends $1014$\,s
  over three seeds and returns no routing, while \dSABRE{}'s $48.5$\,s spans
  three \texttt{SabreLayout} seeds, none of which now aborts under the
  capacity invariant of Section~\ref{sec:relief}.  \texttt{pytket-dqc} is not
  reported on that series: its counts there were measured on an earlier
  network construction and have not been re-run, and the comparison of
  Section~\ref{sec:pytket} is made on the three fixed-architecture suites.
  The 64-qubit block was
  measured in isolation; the scalability block comes from the benchmark
  sweep.}
\label{tab:timing}
\setlength{\tabcolsep}{5pt}
\small
\resizebox{\ifdim\width>\columnwidth \columnwidth\else\width\fi}{!}{%
\begin{tabular}{@{}l r rrr@{}}
\toprule
 & & \TeleSABRE{} & \dSABRE{} & \texttt{pytket-dqc} \\
Circuit & CX & ($1$ seed) & ($3$ seeds) & ($5$ seeds) \\
\midrule
GHZ                     &    63 & 0.04   & 0.05    & 0.3 \\
Graphstate              &    64 & 2.74   & 0.14    & 0.6 \\
AE                      &  1962 & 1.78   & 2.53    & 291.3 \\
QFT                     &  1966 & 10.95  & 3.23    & 490.9 \\
Random$^{\ast\dagger}$  &  1627 & ---          & 6.06    & 171.0 \\
QPEexact                &  2139 & 1.48   & 2.03    & 458.3 \\
QAOA$^{\ddagger}$       &  3920 & 52.32  & 5.29    & 899.4 \\
QNN$^{\dagger}$         &  8126 & 9.87   & 10.06   & 1126.9 \\
Multiplier$^{\dagger}$  & 13040 & 136.25 & 23.70   & 1392.5 \\
\cmidrule(lr){1-5}
\textbf{gmean} (8)      &       & \textbf{4.66} & \textbf{1.87}
                                & \textbf{106.6} \\
\midrule
\multicolumn{5}{@{}l}{\emph{QFT scalability, $5{\times}5$ cores, core count $6\to12\to20$}} \\
\addlinespace[1pt]
QFT 100q          & 3420  & 3.6  & 3.4   & --- \\
QFT 200q          & 7220  & 60.9 & 13.5  & --- \\
QFT 360q$^{\ast}$ & 13300 & ---         & 48.5 & --- \\
\bottomrule
\end{tabular}%
}
\end{table}

\subsection{Benchmark provenance and per-seed counts}
\label{app:provenance}
Table~\ref{tab:circuits} records the CX count and depth of every benchmark circuit after transpilation, and Table~\ref{tab:seeds} the per-seed EPR counts behind every best-of-three figure reported in the main paper.

\begin{table}[t]
\centering
\caption{Benchmark circuit properties (measurements/barriers stripped)}
\label{tab:circuits}
\setlength{\tabcolsep}{5pt}
\begin{tabular}{@{}llrrrr@{}}
\toprule
Suite & Circuit    & Qubits & CX gates & Depth & CX/qubit \\
\midrule
\multirow{6}{*}{25q}
 & AE         & 25 &  558 & 395 & 22.3 \\
 & QFT        & 25 &  580 & 173 & 23.2 \\
 & QNN        & 25 & 1223 & 259 & 48.9 \\
 & Random     & 25 & 1124 & 589 & 45.0 \\
 & GHZ        & 25 &   24 &  27 &  1.0 \\
 & Graphstate & 25 &   25 &  19 &  1.0 \\
\midrule
\multirow{6}{*}{36q}
 & BV         & 36 &   17 &  23 &  0.5 \\
 & DJ         & 36 &   35 &  41 &  1.0 \\
 & W-state    & 36 &   70 & 145 &  1.9 \\
 & VQE-SU2   & 36 &  105 &  56 &  2.9 \\
 & QPEexact   & 36 & 1019 & 347 & 28.3 \\
 & QAOA       & 36 & 1200 & 256 & 33.3 \\
\midrule
\multirow{9}{*}{64q}
 & AE         & 64 & 1962 & 1058 & 30.7 \\
 & QFT        & 64 & 1966 &  446 & 30.7 \\
 & QNN        & 64 & 8126 &  650 & 127.0 \\
 & Random     & 64 & 1627 &  403 & 25.4 \\
 & GHZ        & 64 &   63 &   66 &  1.0 \\
 & Graphstate & 64 &   64 &   25 &  1.0 \\
 & QPEexact   & 64 & 2139 & 626 & 33.4 \\
 & QAOA       & 64 & 3920 & 493 & 61.2 \\
 & Multiplier & 64 & 13040 & 22469 & 203.8 \\
\bottomrule
\end{tabular}
\end{table}

\begin{table}[t]
\centering
\caption{\dSABRE{} (dSE) EPR count under each \texttt{SabreLayout}
  seed, with the best---which is what the main tables report---and the
  median.  \TeleSABRE{} is a single column because its three seeds
  coincide.  Geometric means are over the circuits \TeleSABRE{} completes,
  matching Table~\ref{tab:main}'s basis; the 360-qubit suite has none and
  so has no mean.}
\label{tab:seeds}
\small
\setlength{\tabcolsep}{4pt}
\resizebox{\ifdim\width>\columnwidth \columnwidth\else\width\fi}{!}{%
\begin{tabular}{@{}l rrr rr r@{}}
\toprule
 & \multicolumn{3}{c}{\dSABRE{} by seed} & & & \\
\cmidrule(lr){2-4}
Circuit & 0 & 1 & 2 & best & median & \TeleSABRE{} \\
\midrule
% BODY GENERATED BY code/gen_seed_table.py -- regenerate and paste
\multicolumn{7}{@{}l}{\emph{25 logical qubits, B-grid}} \\
\addlinespace[1pt]
AE & 23 & 23 & 23 & 23 & 23 & 26 \\
GHZ & 1 & 1 & 1 & 1 & 1 & 4 \\
Graphstate & 2 & 2 & 2 & 2 & 2 & 13 \\
QFT & 33 & 33 & 33 & 33 & 33 & 46 \\
QNN & 47 & 48 & 51 & 47 & 48 & 54 \\
Random & 181 & 188 & 197 & 181 & 188 & 271 \\
\cmidrule(lr){1-7}
\textbf{gmean} (6) & & & & \textbf{15.3} & 15.5 & 31.1 \\
\midrule
\multicolumn{7}{@{}l}{\emph{36 logical qubits, B-grid}} \\
\addlinespace[1pt]
BV & 1 & 1 & 2 & 1 & 1 & 6 \\
DJ & 3 & 3 & 3 & 3 & 3 & 4 \\
QAOA & 134 & 136 & 139 & 134 & 136 & 268 \\
QPEexact & 63 & 63 & 62 & 62 & 63 & 190 \\
VQE-SU2 & 12 & 9 & 16 & 9 & 12 & 28 \\
W-state & 7 & 6 & 6 & 6 & 6 & 13 \\
\cmidrule(lr){1-7}
\textbf{gmean} (6) & & & & \textbf{10.5} & 11.1 & 27.6 \\
\midrule
\multicolumn{7}{@{}l}{\emph{64 logical qubits, H-grid}} \\
\addlinespace[1pt]
AE & 165 & 196 & 212 & 165 & 196 & 353 \\
GHZ & 6 & 6 & 6 & 6 & 6 & 16 \\
Graphstate & 20 & 26 & 15 & 15 & 20 & 62 \\
QFT & 180 & 203 & 243 & 180 & 203 & 312 \\
QNN & 659 & 539 & 607 & 539 & 607 & 889 \\
Random & 753 & 763 & 732 & 732 & 753 & --- \\
QPEexact & 220 & 187 & 197 & 187 & 197 & 289 \\
QAOA & 559 & 603 & 555 & 555 & 559 & 956 \\
Multiplier & 1285 & 1529 & 1566 & 1285 & 1529 & 1601 \\
\cmidrule(lr){1-7}
\textbf{gmean} (8) & & & & \textbf{144.7} & 162.6 & 284.5 \\
\midrule
\multicolumn{7}{@{}l}{\emph{100 logical qubits, H-grid $2{\times}3$ of $5{\times}5$}} \\
\addlinespace[1pt]
QFT & 201 & 264 & 262 & 201 & 262 & 248 \\
QPEexact & 246 & 317 & 411 & 246 & 317 & 586 \\
\cmidrule(lr){1-7}
\textbf{gmean} (2) & & & & \textbf{222.4} & 288.2 & 381.2 \\
\midrule
\multicolumn{7}{@{}l}{\emph{200 logical qubits, H-grid $3{\times}4$ of $5{\times}5$}} \\
\addlinespace[1pt]
QFT & 913 & 820 & 766 & 766 & 820 & 1232 \\
QPEexact & 896 & 1042 & 903 & 896 & 903 & --- \\
\cmidrule(lr){1-7}
\textbf{gmean} (1) & & & & \textbf{766.0} & 820.0 & 1232.0 \\
\midrule
\multicolumn{7}{@{}l}{\emph{360 logical qubits, H-grid $4{\times}5$ of $5{\times}5$}} \\
\addlinespace[1pt]
QFT & 2229 & 1890 & 2096 & 1890 & 2096 & --- \\
QPEexact & 3168 & 2709 & 1938 & 1938 & 2709 & --- \\
\bottomrule
\end{tabular}%
}
\end{table}

% -----------------------------------------------------------------------------
\section{Cross-Model Comparisons}
\label{app:crossmodel}

\subsection{pytket-dqc capacity audit}
\label{app:pytketdqc}
\texttt{pytket-dqc}~\cite{pytketdqc2024} uses static hypergraph
partitioning and amortised gate teleportation, while \dSABRE{} uses
dynamic teledata and intra-core routing.
Their EPR units are therefore comparable only after stating the capacity and
entanglement-lifetime assumptions.  Table~\ref{tab:appendix-pytket-capacity}
reports the smallest uniform link-qubit capacity found sufficient to
materialise each \texttt{pytket-dqc} distribution.  A dagger means the
requirement provably exceeds the two ports per B-grid core or the two to three
ports per H-grid core.  Values prefixed by $>$ are lower bounds because the
search budget expired.

\begin{table}[t]
\centering
\caption{\texttt{pytket-dqc} e-bits and required uniform communication
capacity, compared with \dSABRE{} EPR.}
\label{tab:appendix-pytket-capacity}
\small
\begin{tabular}{@{}lrrr@{}}
\toprule
Circuit & e-bits & \dSABRE{} EPR & ports needed \\
\midrule
\multicolumn{4}{@{}l}{\emph{25q B-grid (2 ports/core)}}\\
AE$^\dagger$ & 48 & 23 & 20 \\
GHZ & 3 & 1 & 1 \\
Graphstate & 4 & 2 & 1 \\
QFT$^\dagger$ & 59 & 33 & 7 \\
QNN$^\dagger$ & 114 & 47 & $>12$ \\
Random$^\dagger$ & 468 & 181 & $>7$ \\
\midrule
\multicolumn{4}{@{}l}{\emph{36q B-grid (2 ports/core)}}\\
BV & 3 & 1 & 1 \\
DJ & 3 & 3 & 1 \\
QAOA$^\dagger$ & 150 & 134 & $>13$ \\
QPEexact$^\dagger$ & 45 & 62 & 7 \\
VQE-SU2 & 9 & 9 & 1 \\
W-state & 6 & 6 & 1 \\
\midrule
\multicolumn{4}{@{}l}{\emph{64q H-grid (2--3 ports/core)}}\\
AE$^\dagger$ & 120 & 165 & $>28$ \\
GHZ & 7 & 6 & 1 \\
Graphstate & 6 & 15 & 1 \\
QFT$^\dagger$ & 157 & 180 & $>6$ \\
QNN$^\dagger$ & 530 & 539 & $>6$ \\
Random$^\dagger$ & 895 & 732 & $>4$ \\
QPEexact$^\dagger$ & 75 & 187 & $>5$ \\
QAOA$^\dagger$ & 462 & 555 & $>10$ \\
Multiplier & 2999 & 1285 & $>1$ \\
\bottomrule
\end{tabular}
\end{table}

The table is an executability audit, not a claim that gate teleportation is
inferior in general.  Under its more permissive default model,
\texttt{pytket-dqc} obtains 34.9, 12.2, and 150.2 geometric-mean e-bits on
the 25q, 36q, and 64q suites, respectively, versus \dSABRE{}'s 15.3, 10.5,
and 173.3 EPR pairs.  Section~IV-E explains why these figures must not be
read as a like-for-like win/loss comparison.

\subsection{Entanglement lifetime}
\label{app:dmax}
Every e-bit count above assumes an unbounded entanglement lifetime.
Table~\ref{tab:dmax} re-scores the same distributions with a counter that
caps gates per EPR pair at $D_{\max}$~\cite{bandini2026optimized} and
reproduces \texttt{pytket-dqc}'s own cost exactly at $D_{\max}{=}\infty$.  The 64-qubit comparison crosses
parity between $D_{\max}{=}\infty$ and $10$ and becomes a $59.6\%$
\dSABRE{} advantage at $D_{\max}{=}3$.  \dSABRE{} needs no adjustment: a
teledata pair is consumed by the teleport that creates it.

\begin{table}[t]
\centering
\caption{{Entanglement-lifetime sensitivity on the 64q suite:
  \texttt{pytket-dqc} e-bits when one EPR pair may serve at most
  $D_{\max}$ gates (counter as described in the text).  The
  distributions scored are the ones Table~\ref{tab:fair} reports, loaded
  from disk rather than searched afresh, so the $D_{\max}{=}\infty$
  column \emph{is} that table's \texttt{pytket-dqc} column and the
  counter reproduces its cost exactly, circuit by circuit.  \dSABRE{}
  EPR (Table~\ref{tab:main}) is $D_{\max}$-invariant.
  $\Delta_{3}$ compares \dSABRE{} against the $D_{\max}{=}3$ column;
  negative favours \dSABRE{}.}}
\label{tab:dmax}
\setlength{\tabcolsep}{3pt}
\small
\resizebox{\columnwidth}{!}{%
\begin{tabular}{@{}l r rrrrr r@{}}
\toprule
 & \dSABRE & \multicolumn{5}{c}{\texttt{pytket-dqc} e-bits at $D_{\max}$} & \\
\cmidrule(lr){3-7}
Circuit & EPR & $\infty$ & $10$ & $5$ & $3$ & $1$ & $\Delta_{3}$ (\%) \\
\midrule
% BODY GENERATED BY code/gen_dmax_table.py -- regenerate and paste
ae & 165 & 120 & 219 & 382 & 595 & 1670 & $-72.3$ \\
ghz & 6 & 7 & 7 & 7 & 7 & 7 & $-14.3$ \\
graphstate & 15 & 6 & 6 & 6 & 6 & 6 & $+150.0$ \\
qft & 180 & 157 & 268 & 432 & 652 & 1777 & $-72.4$ \\
qnn & 539 & 530 & 1446 & 2527 & 3891 & 11232 & $-86.1$ \\
random & 732 & 895 & 895 & 939 & 1034 & 1806 & $-29.2$ \\
qpeexact & 187 & 75 & 158 & 288 & 451 & 1282 & $-58.5$ \\
qaoa & 555 & 462 & 832 & 1488 & 2306 & 6464 & $-75.9$ \\
multiplier & 1285 & 2999 & 3869 & 5149 & 7164 & 17882 & $-82.1$ \\
\midrule
\textbf{gmean} $\Delta$ (\%) & & $+15.3$ & $-23.8$ & $-46.0$ & $\mathbf{-59.6}$ & $-80.7$ & \\
\bottomrule
\end{tabular}%
}
\end{table}

\subsection{DMapS comparison}
\label{app:dmaps}
DMapS~\cite{luo2025dmaps} uses a different primitive: its cross-chip movement
is principally a remote SWAP, which exchanges two states and therefore has the
cost of two one-way teledata moves under the EPR accounting used here.  In the nine-circuit
64-qubit suite, \dSABRE{} uses $1.43\times$ fewer EPR pairs in geometric mean;
DMapS is lower only on GHZ and Graphstate, where partition structure is
especially favourable.  This result is reported as a cross-model comparison,
not as a direct baseline: the two routers differ in communication primitive,
occupancy behavior, and initial partitioning.  Table~\ref{tab:dmaps_vs_dsabre} gives the per-circuit head-to-head the main
paper points to; the released artefact contains the machine-readable results
and exact environment instructions.

\begin{table*}[t]
\centering
\caption{EPR and local-SWAP head-to-head: DMapS vs.\ \dSABRE{}.
  DMapS is best of five KaHyPar seeds (its own protocol); \dSABRE{} is
  best of three SabreLayout seeds, matching the main-text headline
  protocol.  ``rCX'' and ``rSWAP'' are the cross-chip CX and
  cross-chip SWAP counts in DMapS's output circuit;
  $\textsf{DMapS EPR} = \textsf{rCX} + 2\!\cdot\!\textsf{rSWAP}$.
  $\textsf{EPR}\times = \textsf{DMapS EPR} / \dSABRE{}\;\textsf{EPR}$;
  $\textsf{lSW}\times$ analogously for local SWAPs.
  Values $<1$ mean DMapS uses fewer EPRs (resp.\ local SWAPs); values
  $>1$ mean \dSABRE{} wins.  Rows marked with $\dagger$ are circuits that
  only run after our upstream DMapS patches~(3)/(4).}
\label{tab:dmaps_vs_dsabre}
\setlength{\tabcolsep}{4pt}
\small
\begin{tabular}{@{}l l r r r r r r r r r@{}}
\toprule
Suite & Circuit & CX & \dSABRE{} EPR & \dSABRE{} lSW & rCX & rSWAP & DMapS EPR & DMapS lSW & $\textsf{EPR}\times$ & $\textsf{lSW}\times$ \\
\midrule
64q & ghz        &   63 &    6 &    39 &    3 &    0 &      3 &      17 & $0.50$ & $0.44$ \\
64q & graphstate &   64 &   15 &    88 &    4 &    0 &      4 &      24 & $0.27$ & $0.27$ \\
64q & ae$^{\dagger}$ & 1962 &  165 &  1325 &  116 &  115 &    346 &    1119 & \textbf{$2.10$} & $0.84$ \\
64q & qft        & 1966 &  180 &  1480 &  148 &  122 &    392 &    1152 & \textbf{$2.18$} & $0.78$ \\
64q & qpeexact   & 2139 &  187 &  1576 &  158 &  125 &    408 &    1136 & \textbf{$2.18$} & $0.72$ \\
64q & random     & 1627 &  732 &  3603 &  378 &  380 &   1138 &    3446 & \textbf{$1.55$} & $0.96$ \\
64q & qaoa       & 3920 &  555 &  3734 &  456 &  400 &   1256 &    3199 & \textbf{$2.26$} & $0.86$ \\
64q & qnn        & 8126 &  539 &  5260 &  405 &  264 &    933 &    3762 & \textbf{$1.73$} & $0.72$ \\
64q & multiplier & 13040 & 1285 &  8969 & 2376 &  784 &   3944 &    8514 & \textbf{$3.07$} & $0.95$ \\
\cmidrule(lr){1-11}
64q & \textbf{gmean (9)} & & \textbf{173.3} & \textbf{1236} & & & \textbf{247.6} & \textbf{841} & \textbf{$1.43$} & $0.68$ \\
\bottomrule
\end{tabular}
\end{table*}
